% !TeX program = pdflatex
% Projektbericht MiniC++ -- Entwurf v0.1 (aus docs/report.py übertragen)
\documentclass[11pt,a4paper,parskip=half,headsepline,toc=flat]{scrartcl}

% ------------------------------------------------------------------ Grundlagen
\usepackage[utf8]{inputenc}
\usepackage[T1]{fontenc}
\usepackage[ngerman]{babel}
\usepackage[autostyle,german=quotes]{csquotes}
\usepackage{lmodern}
\usepackage[scaled=0.85]{beramono}
\usepackage{textcomp}
\usepackage{microtype}
\usepackage[a4paper,left=2.3cm,right=2.3cm,top=2.4cm,bottom=2.3cm,headsep=0.6cm]{geometry}

\usepackage[table]{xcolor}
\usepackage{array,booktabs,tabularx}
\usepackage{enumitem}
\usepackage{caption}
\usepackage{listings}
\usepackage[most]{tcolorbox}
\usepackage{tikz}
\usetikzlibrary{arrows.meta,positioning,calc,decorations.pathreplacing}
\usepackage{pgfplots}
\pgfplotsset{compat=1.18}
\usepackage[automark]{scrlayer-scrpage}
\usepackage[hidelinks,pdfusetitle]{hyperref}

\hypersetup{
  pdftitle={MiniC++ Interpreter-Ökosystem -- Projektbericht (Entwurf)},
  pdfauthor={Clemens Vogtländer, Dennis Gorpinic},
  pdfsubject={Technische Dokumentation}
}

% Das "-Kürzel von ngerman stört bei Code wie {"error": ...}; deutsche Anführungszeichen via \enquote.
\AtBeginDocument{\shorthandoff{"}}
\setcounter{secnumdepth}{2}
\setcounter{tocdepth}{2}
\emergencystretch=2em

% ------------------------------------------------------------------ Farben
\definecolor{navy}{HTML}{1B3556}
\definecolor{accent}{HTML}{2A78D6}
\definecolor{ink}{HTML}{1A1A1A}
\definecolor{ink2}{HTML}{52514E}
\definecolor{rule}{HTML}{D7D6D1}
\definecolor{codebg}{HTML}{F4F5F7}
\definecolor{headbg}{HTML}{E8EEF6}
\definecolor{zebra}{HTML}{F8F8F6}
\definecolor{todobg}{HTML}{FFF5D6}
\definecolor{todobd}{HTML}{C98500}
\definecolor{notebg}{HTML}{EAF2FB}
\definecolor{lexbg}{HTML}{DFEAF7}
\definecolor{gridc}{HTML}{E4E3DF}
\definecolor{seqA}{HTML}{9EC5F0}
\definecolor{seqB}{HTML}{2A78D6}
\definecolor{seqC}{HTML}{154A8A}
\definecolor{orange}{HTML}{EB6834}

% ------------------------------------------------------------------ Überschriften
\addtokomafont{disposition}{\color{navy}}
\addtokomafont{section}{\Large}
\addtokomafont{subsubsection}{\color{ink}}

% ------------------------------------------------------------------ Kopf- und Fußzeile
\pagestyle{scrheadings}
\clearpairofpagestyles
\ihead{\footnotesize\color{ink2}MiniC++ -- Interpreter-Ökosystem \textperiodcentered{} Technische Dokumentation}
\ohead{\footnotesize\color{todobd}ENTWURF v0.1 \textperiodcentered{} 27.09.2026}
\cfoot*{\footnotesize\color{ink2}\pagemark}
\setkomafont{pageheadfoot}{\normalfont}

% ------------------------------------------------------------------ Hilfsmakros
\newcommand{\code}[1]{\texttt{#1}}
\newcommand{\bs}{\textbackslash}
\newcommand{\dq}{\textquotedbl}
\newcommand{\sq}{\textquotesingle}

% Tabellen: Kopfzeile mit Hintergrund, Zebrastreifen
\newcolumntype{L}[1]{>{\raggedright\arraybackslash}p{#1\linewidth}}
\newcommand{\thead}{\rowcolor{headbg}}
\newcommand{\hcell}[1]{\textbf{\color{navy}#1}}
\captionsetup{font={small,it},labelfont={small,it},textfont={color=ink2},labelsep=colon}
\captionsetup[table]{position=bottom}
\newenvironment{reporttable}[1][]%
  {\begin{table}[htbp]\centering\small\renewcommand{\arraystretch}{1.25}\rowcolors{2}{white}{zebra}}%
  {\end{table}}

% Listings
\lstset{
  basicstyle=\ttfamily\footnotesize,
  backgroundcolor=\color{codebg},
  frame=leftline, framerule=2pt, rulecolor=\color{accent},
  xleftmargin=8pt, framexleftmargin=6pt, framesep=6pt,
  aboveskip=8pt, belowskip=4pt,
  columns=fullflexible, keepspaces=true, showstringspaces=false,
  breaklines=true,
  literate={ä}{{\"a}}1 {ö}{{\"o}}1 {ü}{{\"u}}1 {Ä}{{\"A}}1 {Ö}{{\"O}}1 {Ü}{{\"U}}1 {ß}{{\ss}}1
}
\renewcommand{\lstlistingname}{Listing}

% TODO- und Hinweis-Kästen
\newtcolorbox{todo}{enhanced,breakable,colback=todobg,colframe=todobd,boxrule=0pt,leftrule=2.5pt,
  arc=0pt,outer arc=0pt,left=8pt,right=8pt,top=5pt,bottom=5pt,fontupper=\small,
  before upper={\textbf{TODO (Entwurf):}\ }}
\newtcolorbox{hinweis}{enhanced,breakable,colback=notebg,colframe=accent,boxrule=0pt,leftrule=2.5pt,
  arc=0pt,outer arc=0pt,left=8pt,right=8pt,top=5pt,bottom=5pt,fontupper=\small,
  before upper={\textbf{Hinweis:}\ }}

\setlist[itemize]{leftmargin=1.4em,itemsep=2pt,topsep=2pt}

% TikZ-Stile für die Diagramme
\tikzset{
  >=Stealth,
  dbox/.style={draw=navy, fill=headbg, rounded corners=2pt, align=center, inner sep=3pt,
               line width=0.6pt, font=\footnotesize},
  sub/.style={font=\scriptsize\itshape, color=ink2},
  lbl/.style={font=\scriptsize\itshape, color=ink2},
  arr/.style={->, color=ink2, line width=0.7pt},
}
\newcommand{\boxtext}[2]{\textbf{#1}\\[-1pt]{\scriptsize\itshape\color{ink2}#2}}

% ==================================================================
\begin{document}

% ------------------------------------------------------------------ Titelseite
\begin{titlepage}
\begin{tikzpicture}[remember picture,overlay]
  \fill[navy] (current page.north west) rectangle ([yshift=-9.5cm]current page.north east);
  \fill[accent] ([yshift=-9.5cm]current page.north west) rectangle ([yshift=-9.65cm]current page.north east);
\end{tikzpicture}
\vspace*{0.4cm}
{\color{white}\sffamily\bfseries\fontsize{24}{29}\selectfont
Entwurf und Implementierung eines interaktiven Interpreter-Ökosystems für eine eingeschränkte C++-Teilmenge\par}
\vspace{8pt}
{\color[HTML]{CFE0F5}\sffamily\large mit Typprüfung, REPL, Language-Server und GenAI-Assistenz\par}
\vspace{3.6cm}

{\color{navy}\sffamily\bfseries\Large Projektbericht \textperiodcentered{} Technische Dokumentation\par}
\vspace{10pt}
{\small\renewcommand{\arraystretch}{1.5}
\begin{tabularx}{\linewidth}{@{}>{\bfseries\color{navy}}p{0.27\linewidth}X@{}}
Autoren               & Clemens Vogtländer, Dennis Gorpinic \\ \arrayrulecolor{rule}\hline
Lehrveranstaltung     & Compilerbau (CB) -- Erweiterung von Aufgabenblatt 08 \\ \hline
Hochschule / Semester & [TODO: Hochschule, Studiengang, Semester, Betreuer/in] \\ \hline
Repository            & cpp-subset-interpreter, Branch interpreter (Stand 27.09.2026) \\ \hline
Dokumentstatus        & Entwurf v0.1 -- nicht zur Abgabe bestimmt \\ \hline
\end{tabularx}}
\vspace{1cm}

{\sffamily\bfseries Kurzfassung\par}
\vspace{2pt}
Dieser Bericht dokumentiert Entwurf und Umsetzung von \emph{MiniC++}, einem formal abgegrenzten Subdialekt von C++,
sowie des darum gebauten Werkzeug-Ökosystems. Kern ist ein in Java~21 implementierter Tree-Walking-Interpreter mit
ANTLR4-Frontend, zweiphasigem Resolver, deklarationsbasierter Typprüfung mit Überladungsauflösung und einem
Laufzeitmodell für Referenzen, Zeiger, Heap-Objekte, Slicing und virtuellen Dispatch. Darauf setzen eine REPL, ein
Language-Server nach LSP~3.17 mit VS-Code-Extension und ein REST-basierter GenAI-Assistenzserver (Ollama) auf, dessen
Antworten durch den echten Compiler abgesichert werden. 218 automatisierte Tests, ein differenzieller Vergleich mit
GCC in der CI und JMH-Benchmarks bilden die Evaluationsgrundlage.
\end{titlepage}

% ------------------------------------------------------------------ Inhaltsverzeichnis
\tableofcontents
\clearpage

% ================================================================== 1 Einleitung
\section{Einleitung}
\subsection{Ausgangslage und Zielsetzung}
Aufgabenblatt 08 der Lehrveranstaltung Compilerbau verlangt einen Interpreter für eine Teilmenge von C++ mit Klassen,
Einfachvererbung, Referenzen und virtuellen Methoden. Das hier beschriebene Projekt nimmt diese Aufgabe als Kern und
erweitert sie zu einem vollständigen \emph{Ökosystem}: Neben der reinen Ausführung von Programmen soll die Sprache so
benutzbar sein, wie man es von einer produktiven Programmiersprache erwartet -- mit einer interaktiven Shell,
Editor-Unterstützung und KI-gestützter Code-Assistenz.

Daraus ergeben sich drei übergeordnete Ziele:
\begin{itemize}
  \item \textbf{Korrektheit gegenüber C++.} Jedes gültige MiniC++-Programm soll sich genau so verhalten wie das
    entsprechende Standard-C++-Programm. Wo C++ undefiniertes Verhalten zulässt (hängende Zeiger, doppeltes
    \code{delete}), soll MiniC++ stattdessen einen präzisen Laufzeitfehler melden.
  \item \textbf{Wiederverwendbarkeit des Frontends.} Lexer, Parser, Resolver und Typprüfung sollen nicht nur dem
    Interpreter dienen, sondern über eine schmale Fassade auch dem Language-Server und dem GenAI-Server.
  \item \textbf{Messbarkeit.} Korrektheit, Performance und Skalierbarkeit sollen mit reproduzierbaren Werkzeugen
    (JUnit, GCC-Vergleich, JMH) belegt werden.
\end{itemize}

\subsection{Projektumfang}
Das System gliedert sich in vier Säulen, die in diesem Bericht jeweils ein eigenes Kapitel erhalten:

\begin{reporttable}
\begin{tabular}{L{0.18}L{0.52}L{0.22}}
\thead \hcell{Säule} & \hcell{Inhalt} & \hcell{Verantwortung} \\ \midrule
Kern-Interpreter & Lexer, Parser, AST, Resolver, Typprüfung, Tree-Walking-Interpreter, REPL, CLI, C++-Export & gemeinsam \\
LSP-Server & Language Server Protocol 3.17 über LSP4J, VS-Code-Extension & Clemens Vogtländer \\
GenAI-Server & REST-Endpunkte für Vervollständigung, Erklärung, Refactoring und Fehlersuche mit Ollama & Dennis Gorpinic \\
Evaluation & Unit- und Golden-File-Tests, differenzieller GCC-Vergleich, JMH-Benchmarks & gemeinsam \\ \bottomrule
\end{tabular}
\caption{Die vier Säulen des Projekts}
\end{reporttable}

\subsection{Aufbau des Dokuments}
Kapitel~\ref{sec:sprache} definiert den Sprachumfang von MiniC++ einschließlich aller Präzisierungen, die über die
Aufgabenstellung hinaus getroffen wurden. Kapitel~\ref{sec:architektur} beschreibt die Gesamtarchitektur und die
Architekturentscheidungen. Kapitel~\ref{sec:kern} behandelt den Kern-Interpreter Phase für Phase,
Kapitel~\ref{sec:lsp} den Language-Server und Kapitel~\ref{sec:genai} den GenAI-Assistenzserver.
Kapitel~\ref{sec:evaluation} fasst Teststrategie und Messergebnisse zusammen, Kapitel~\ref{sec:organisation} die
Projektorganisation. Kapitel~\ref{sec:ausblick} diskutiert Grenzen und mögliche Weiterentwicklungen. Der Anhang
enthält eine Kurzreferenz der Kommandozeile und die Kernregeln der Grammatik.

Konventionen: Bezeichner aus dem Quellcode sind in \code{Festbreitenschrift} gesetzt; Pfade beziehen sich auf die
Wurzel des Repositorys. Das Java-Basispaket \code{de.cvogtlaender} wird in Klassennamen weggelassen.

% ================================================================== 2 Sprache
\section{Die Sprache MiniC++}\label{sec:sprache}
MiniC++ ist eine echte Teilmenge von C++17: Jedes gültige MiniC++-Programm ist -- nach einer rein mechanischen
Übersetzung, siehe Abschnitt~\ref{sec:cppexport} -- ein gültiges C++-Programm mit identischer Ausgabe. Dieses Kapitel
beschreibt den Sprachumfang aus Sicht des Programmierers.

\subsection{Typen, Variablen und Werte}
MiniC++ kennt fünf Basistypen, benutzerdefinierte Klassentypen sowie daraus gebildete Referenz- und Zeigertypen.
Variablen werden mit \code{T x;} oder \code{T x = expr;} deklariert; für Klassentypen ist \code{T x(a, b);} eine
Kurzform von \code{T x = T(a, b);}. Globale Variablen gibt es nicht.

\begin{reporttable}
\begin{tabular}{L{0.13}L{0.5}L{0.29}}
\thead \hcell{Typ} & \hcell{Wertebereich / Semantik} & \hcell{Standardwert ohne Initialisierer} \\ \midrule
\code{int} & 32 Bit, Zweierkomplement mit Überlauf (wie \code{g++ -fwrapv}) & \code{0} \\
\code{bool} & \code{true} / \code{false} & \code{false} \\
\code{char} & Zeichen inkl.\ Escape-Sequenzen \code{\bs n \bs t \bs\sq{} \bs\bs{} \bs 0} u.\,a. & \code{\sq\bs 0\sq} \\
\code{string} & Zeichenkette mit Wertsemantik & \code{\dq\dq} \\
\code{void} & nur als Rückgabetyp & -- \\
Klasse \code{A} & Objekt mit Wertsemantik (Kopie bei Zuweisung) & parameterloser Konstruktor \\
\code{T\&} & Referenz; Initialisierung obligatorisch, keine Neubindung & -- (Fehler) \\
\code{T*}, \code{T**} & Zeiger, auch mehrstufig; keine Arithmetik & \code{nullptr} \\ \bottomrule
\end{tabular}
\caption{Typen von MiniC++ und ihre Standardwerte}
\end{reporttable}

\subsection{Ausdrücke und Operatoren}
Operatorpräzedenz und Assoziativität folgen dem C++-Standard. Die Typregeln sind bewusst strenger als in C++: Es gibt
keine impliziten arithmetischen Konvertierungen, und beide Operanden eines binären Operators müssen denselben Typ
haben.

\begin{reporttable}
\begin{tabular}{L{0.16}L{0.28}L{0.48}}
\thead \hcell{Stufe} & \hcell{Operatoren} & \hcell{zulässige Operandentypen} \\ \midrule
1 (höchste) & Aufruf \code{f()}, Member \code{.} \code{->} & Funktionen, Methoden, Objekte, Zeiger auf Objekte \\
2 & unär \code{! + - * \&}, \code{new} & \code{!} nur \code{bool}; \code{+ -} nur \code{int}; \code{*} Zeiger; \code{\&} L-Werte \\
3 & \code{* / \%} & \code{int} \\
4 & \code{+ -} & \code{int} \\
5 & \code{< <= > >=} & \code{int}, \code{char} \\
6 & \code{== !=} & \code{int}, \code{char}, \code{bool}, \code{string}, kompatible Zeiger \\
7 & \code{\&\&} & \code{bool} (Kurzschlussauswertung) \\
8 & \code{||} & \code{bool} (Kurzschlussauswertung) \\
9 (niedrigste) & \code{=} (rechtsassoziativ) & L-Wert links; Slicing bei Klassen erlaubt \\ \bottomrule
\end{tabular}
\caption{Operatoren nach Präzedenz}
\end{reporttable}

Implizite Konvertierung nach \code{bool} findet ausschließlich in \code{if}- und \code{while}-Bedingungen statt; dort
sind auch \code{int}, \code{char} und Zeiger erlaubt. Operanden und Argumente werden strikt von links nach rechts
ausgewertet. Division und Modulo durch null sind Laufzeitfehler.

\subsection{Anweisungen, Funktionen und Überladung}
Der Kontrollfluss umfasst Blöcke, \code{if}/\code{else}, \code{while} und \code{return}; \code{for}, \code{switch},
\code{break} und \code{continue} gehören nicht zur Sprache. Funktionen können über Name, Arität und Parametertypen
einschließlich der \code{\&}-Markierung überladen werden. Eingebaut sind \code{print\_bool}, \code{print\_int},
\code{print\_char} und \code{print\_string}, die ihren Wert gefolgt von einem Zeilenumbruch ausgeben. Einstiegspunkt
ist \code{int main()} oder \code{void main()}; \code{int main()} ohne \code{return} liefert 0, und der Rückgabewert
wird zum Exit-Code des Prozesses.

\subsection{Klassen, Vererbung und Polymorphie}
Klassen werden als \code{class A \{ public: ... \};} definiert; alle Member sind öffentlich. Fehlt ein Konstruktor,
wird ein parameterloser synthetisiert. Einfachvererbung erfolgt über \code{class D : public B}, wobei vor dem Rumpf
eines Konstruktors von \code{D} implizit der parameterlose Konstruktor von \code{B} läuft. Die Namensauflösung
innerhalb von Methoden folgt der Reihenfolge \emph{lokal $\rightarrow$ eigene Member $\rightarrow$ geerbte Member
$\rightarrow$ global}; wie in C++ verdecken Member einer abgeleiteten Klasse gleichnamige Member der Basisklasse.
\code{this} existiert nicht.

Polymorphie entsteht über \code{virtual}-Methoden, die durch Referenzen oder Zeiger aufgerufen werden. Eine Methode,
die eine virtuelle Methode überschreibt, ist selbst virtuell. Während der Konstruktor einer Basisklasse läuft, wird --
exakt wie in C++ -- deren eigene Version aufgerufen. Die Zuweisung eines abgeleiteten Objekts an eine Variable vom
Basistyp kopiert nur den Basisanteil (\emph{Slicing}).

\begin{lstlisting}[language=C++,caption={Vererbung, virtueller Dispatch und Slicing (gekürzt aus \code{examples/features.cpp})}]
class Account {
public:
  int balance;
  virtual int monthlyFee() { return 5; }
  void endOfMonth() { balance = balance - monthlyFee(); }  // dynamischer Dispatch
};
class StudentAccount : public Account {
public:
  int monthlyFee() { return 0; }                             // implizit virtual
};
int main() {
  StudentAccount bob;
  Account& any = bob;      print_int(any.monthlyFee());      // 0  (Polymorphie)
  Account copy = bob;      print_int(copy.monthlyFee());     // 5  (Slicing)
  Account* p = new StudentAccount();
  print_int(p->monthlyFee());                                // 0
  delete p;
  return 0;
}
\end{lstlisting}

\subsection{Referenzen und Zeiger}
Referenzen (\code{T\& r = x;}, Parameter \code{T\& p}) sind zweite Namen für einen existierenden Speicherort; eine
Zuweisung an \code{r} schreibt in \code{x}. Referenzfelder, Referenz-Rückgaben und nicht initialisierte Referenzen
sind nicht erlaubt. Zeiger (\code{T*}, \code{T**}, auch als Parameter \code{T*\&}, Feld oder Rückgabewert)
unterstützen den Adressoperator \code{\&x}, die Dereferenzierung \code{*p} (lesend und schreibend), Memberzugriff
\code{p->m}, das Literal \code{nullptr}, Heap-Allokation mit \code{new T} bzw.\ \code{new T(args)} und Freigabe mit
\code{delete p;}. Implizit konvertiert werden nur \code{D*} $\rightarrow$ \code{B*} und \code{nullptr}
$\rightarrow$ \code{T*}.

Eine Besonderheit, die direkt aus der C++-Grammatik folgt: Die Anweisung \code{a * b;} ist eine Deklaration eines
Zeigers \code{b} vom Typ \code{a*} und keine Multiplikation.

\subsection{Bewusste Einschränkungen und Präzisierungen}
Nicht unterstützt werden Zeigerarithmetik, \code{void*}, Casts, Arrays, \code{++}/\code{--}, zusammengesetzte
Zuweisungen, \code{break}/\code{continue}, Mehrfachvererbung, Templates, \code{static}, \code{const}, \code{this},
globale Variablen, Initialisierungslisten, Destruktoren und reine Funktionsdeklarationen. Wo die Aufgabenstellung
offen war, wurden folgende Festlegungen getroffen:
\begin{itemize}
  \item Was in C++ undefiniertes Verhalten wäre, ist ein \textbf{Laufzeitfehler}: Dereferenzieren von
    \code{nullptr}, hängende Zeiger und Referenzen auf Variablen, deren Gültigkeitsbereich endete, Zugriff auf
    gelöschte Heap-Objekte, doppeltes \code{delete} sowie \code{delete} auf nicht per \code{new} erzeugte Objekte.
    Nicht freigegebener Speicher ist dagegen kein Fehler.
  \item Mehr als 100\,000 verschachtelte Aufrufe gelten als Stapelüberlauf und werden als Laufzeitfehler gemeldet.
  \item Überladungsauflösung: Zuerst werden exakte Treffer gesucht; nur wenn keiner existiert, werden die
    Konvertierungen Derived $\rightarrow$ Base, \code{D*} $\rightarrow$ \code{B*} und \code{nullptr} $\rightarrow$
    \code{T*} berücksichtigt. Mehr als ein bester Kandidat ist ein Fehler.
  \item Felder dürfen in abgeleiteten Klassen nicht erneut deklariert werden.
  \item Funktionen und Klassen dürfen vor ihrer Definition verwendet werden (\emph{define-after-use}), Variablen
    nicht.
\end{itemize}

% ================================================================== 3 Architektur
\section{Systemarchitektur}\label{sec:architektur}
\subsection{Modulstruktur}
Das Repository ist ein Gradle-Multiprojekt mit vier Java-Modulen und einem TypeScript-Projekt. Alle Werkzeuge hängen
ausschließlich vom Modul \code{interpreter} ab; untereinander gibt es keine Abhängigkeiten. Damit bleibt der Kern frei
von LSP-, HTTP- oder Benchmark-Bibliotheken.

\begin{figure}[htbp]
\centering
\begin{tikzpicture}
  \node[dbox, fill=lexbg, minimum width=6.6cm, minimum height=1.5cm, font=\normalsize] (core) at (0,0)
    {\textbf{interpreter}\\[-1pt]{\scriptsize\itshape\color{ink2}Grammatik \textperiodcentered{} AST \textperiodcentered{} Resolver \textperiodcentered{} Typprüfung \textperiodcentered{} Laufzeit \textperiodcentered{} REPL \textperiodcentered{} CLI}};
  \node[dbox, minimum width=3.4cm, minimum height=1cm] (lsp)   at (-5.6,2.4) {\boxtext{lsp}{LSP4J-Server (stdio)}};
  \node[dbox, minimum width=3.4cm, minimum height=1cm] (bench) at (0,2.4)    {\boxtext{benchmark}{JMH-Messungen}};
  \node[dbox, minimum width=3.4cm, minimum height=1cm] (mcp)   at (5.6,2.4)  {\boxtext{mcp}{GenAI-REST-Server}};
  \foreach \m in {lsp,bench,mcp} \draw[arr] (\m.south) -- (core);
  \node[dbox, fill=white, minimum width=3cm, minimum height=1.5cm] (vsc) at (-6.4,0) {\boxtext{vscode}{TypeScript-Client}};
  \node[dbox, fill=white, draw=ink2, minimum width=3cm, minimum height=1.5cm] (oll) at (6.4,0) {\boxtext{Ollama}{lokales LLM (HTTP)}};
  \draw[arr, dashed] (vsc.north -| lsp.south) -- (lsp.south) node[lbl, midway, right] {startet, JSON-RPC};
  \draw[arr, dashed] (mcp.south) -- (oll.north -| mcp.south) node[lbl, midway, right] {/api/generate};
  \node[dbox, fill=white, draw=rule, font=\scriptsize, minimum width=2.1cm] (e1) at (-2.3,-2.3) {ANTLR 4 Runtime};
  \node[dbox, fill=white, draw=rule, font=\scriptsize, minimum width=2.1cm] (e2) at (0,-2.3)    {JUnit 5};
  \node[dbox, fill=white, draw=rule, font=\scriptsize, minimum width=2.1cm] (e3) at (2.3,-2.3)  {Gradle 9.5 Wrapper};
  \draw[arr] (core.south) -- (e2.north) node[lbl, midway, right] {baut auf};
\end{tikzpicture}
\caption{Module und ihre Abhängigkeiten}
\end{figure}

\begin{reporttable}
\begin{tabular}{L{0.15}L{0.43}L{0.22}L{0.1}}
\thead \hcell{Modul} & \hcell{Aufgabe} & \hcell{Quellzeilen (main / test)} & \hcell{Tests} \\ \midrule
\code{interpreter} & Sprache, Laufzeit, REPL, CLI, C++-Export & 6\,421 / 1\,680 & 175 \\
\code{lsp} & Language-Server (LSP4J 1.0.0) & 2\,295 / 590 & 22 \\
\code{mcp} & GenAI-REST-Server (JDK-\code{HttpServer}, Gson) & 618 / 363 & 21 \\
\code{benchmark} & JMH-Benchmarks (JMH 1.37) & 262 / -- & -- \\
\code{vscode} & VS-Code-Client (TypeScript) & 86 & -- \\ \bottomrule
\end{tabular}
\caption{Module des Repositorys (Stand 27.09.2026, inkl.\ Grammatik und Testprogramme)}
\end{reporttable}

\subsection{Verarbeitungspipeline}
Die Verarbeitung eines Programms folgt der klassischen Phasenstruktur eines Compilers. Jede Phase ist als eigene
Klasse realisiert; ab dem AST sind alle Phasen Implementierungen der generischen Schnittstelle
\code{AstVisitor<T>}. Eine Phase läuft nur, wenn die vorherige fehlerfrei war -- so kann sich jede Phase auf die
Invarianten der vorherigen verlassen (etwa: der Interpreter findet jeden Bezeichner bereits aufgelöst und jeden
Ausdruck typisiert vor).

\begin{figure}[htbp]
\centering
\begin{tikzpicture}[stage/.style={dbox, minimum width=2.15cm, text width=2.05cm, minimum height=1.15cm, font=\footnotesize}]
  \foreach \t/\s/\c [count=\i from 0] in
     {Quelltext/String/white, Lexer/MiniCppLexer/lexbg, Parser/MiniCppParser/lexbg,
      AST-Aufbau/ASTBuildVisitor/headbg, Resolver/ASTResolveVisitor/headbg,
      Typprüfung/TypeCheckVisitor/headbg, Interpreter/Interpreter/headbg}
    \node[stage, fill=\c] (s\i) at (\i*2.36,0) {\textbf{\t}\\[-1pt]{\tiny\itshape\color{ink2}\s}};
  \foreach \i [evaluate=\i as \j using int(\i+1)] in {0,...,5}
    \draw[arr] (s\i) -- (s\j);
  \foreach \a/\b/\t in {1/2/Token-Strom, 2/3/Parse-Tree, 3/4/AST + Quellbereiche, 4/5/gebundene Namen, 5/6/annotierte Typen}
    \node[lbl, text width=2.1cm, align=center, anchor=south] at ($(s\a.north)!0.5!(s\b.north)+(0,0.08)$) {\t};
  \draw[decorate, decoration={brace, mirror, amplitude=4pt}, accent, thick]
    ([yshift=-3pt]s1.south west) -- ([yshift=-3pt]s2.south east)
    node[midway, below=5pt, font=\scriptsize\itshape, text=accent, align=center] {generiert aus MiniCpp.g4\\(ANTLR 4.13.2)};
  \draw[decorate, decoration={brace, mirror, amplitude=4pt}, accent, thick]
    ([yshift=-3pt]s3.south west) -- ([yshift=-3pt]s5.south east)
    node[midway, below=5pt, font=\scriptsize\itshape, text=accent] {statische Analyse $\rightarrow$ \code{MiniCpp.compile()}};
  \draw[decorate, decoration={brace, mirror, amplitude=4pt}, accent, thick]
    ([yshift=-3pt]s6.south west) -- ([yshift=-3pt]s6.south east)
    node[midway, below=5pt, font=\scriptsize\itshape, text=accent] {\code{MiniCpp.execute()}};
  \node[font=\scriptsize, color=ink2, anchor=west] at (-1.05,-2.55) {Nutzer der Fassade:};
  \foreach \t/\x in {REPL/6.0, LSP-Server/8.6, GenAI-Server/11.2}
    \node[dbox, fill=white, draw=ink2, minimum width=2.3cm, font=\scriptsize] at (\x,-2.55) {\t};
\end{tikzpicture}
\caption{Verarbeitungspipeline und Fassade \code{MiniCpp}}
\end{figure}

\subsection{Die Fassade MiniCpp}
Alle Nutzer greifen über die Klasse \code{interpreter.MiniCpp} auf die Pipeline zu. Die Ergebnisse sind als
unveränderliche Java-Records modelliert, die neben dem eigentlichen Ergebnis stets eine Liste von Diagnosen tragen:

\begin{lstlisting}[language=Java,caption={Öffentliche Schnittstelle der Fassade \code{MiniCpp} (vereinfacht)}]
record ParseResult<T>(T tree, List<Diagnostic> diagnostics, boolean incomplete)
record Compilation(Program program, GlobalScope globals, List<Diagnostic> diagnostics)
record RunResult(int exitCode, List<Diagnostic> diagnostics)

static ParseResult<Program>   parseProgram(String source)
static ParseResult<ReplInput> parseReplInput(String source)
static Compilation            compile(String source)          // Parsen + Resolver + Typprüfung
static Compilation            analyze(Program p, boolean requireMain)
static RunResult              run(String source, PrintStream out)
static RunResult              execute(Compilation c, PrintStream out)
\end{lstlisting}

Eine \code{Diagnostic} besteht aus der Phase (\code{SYNTAX}, \code{RESOLVE}, \code{TYPE}, \code{RUNTIME}), einem
Quellbereich (Zeilen 1-basiert, Spalten 0-basiert, Ende exklusiv) und einer Meldung. Die Kommandozeile formatiert sie
im von GCC und Clang bekannten Format \code{datei.cpp:3:7: error: use of undeclared identifier \sq y\sq}; der
Language-Server rechnet sie in LSP-Positionen um. Die Trennung von Ergebnis und Darstellung ist der Grund, weshalb
alle drei Werkzeuge dieselben Fehlermeldungen zeigen.

\subsection{Architekturentscheidungen}
Die grundlegenden Entscheidungen wurden zu Projektbeginn gemeinsam getroffen und als Architecture Decision Records
festgehalten. Die folgende Tabelle fasst sie mit ihren Konsequenzen zusammen.

\begin{reporttable}
\begin{tabular}{L{0.23}L{0.38}L{0.31}}
\thead \hcell{Entscheidung} & \hcell{Begründung} & \hcell{Konsequenz im Projekt} \\ \midrule
Java 21 als Implementierungssprache & ausgereifte Standardbibliothek, Klassenhierarchien und Visitor-Pattern passen
  zur AST-Modellierung; Records und Pattern-Matching verkürzen den Code & tiefe Rekursion des Tree-Walkers erfordert
  einen eigenen Thread mit 512 MiB Stack (Abschnitt~\ref{sec:laufzeit}) \\
ANTLR4 als Parsergenerator & De-facto-Standard auf der JVM, Lexer und Parser aus einer \code{.g4}-Datei,
  ALL(*)-Parsing, eingebaute Fehlerbehandlung & Parse-Tree muss in einen eigenen AST übersetzt werden; Grammatik wird
  beim Build generiert \\
Handgeschriebener AST statt Parse-Tree & Phasen sollen auf einer stabilen, typisierten Struktur arbeiten und
  Annotationen (Bindung, Typ) speichern & zusätzliche Klasse \code{ASTBuildVisitor} (757 Zeilen) \\
VS Code als LSP-Zielplattform & gut dokumentierte Extension-API, \code{vscode-languageclient} & Server bleibt
  editorneutral (stdio, LSP4J) \\
Ollama als GenAI-Provider & lokal, ohne API-Schlüssel, Kosten oder Datenschutzbedenken & Provider-Schnittstelle mit
  Mock für deterministische Tests \\ \bottomrule
\end{tabular}
\caption{Architekturentscheidungen}
\end{reporttable}

\begin{reporttable}
\begin{tabular}{L{0.55}L{0.37}}
\thead \hcell{Komponente} & \hcell{Version} \\ \midrule
JDK (Ziel / Entwicklung) & 21 / 25 \\
Gradle (Wrapper) & 9.5.1 \\
ANTLR & 4.13.2 \\
Eclipse LSP4J & 1.0.0 \\
Gson & 2.11.0 \\
JUnit Jupiter & 5.11.4 \\
JMH & 1.37 \\
vscode-languageclient / VS Code Engine & 10.1 / $\geq$ 1.91 \\
CI & GitHub Actions, Temurin 21, Node 22 \\ \bottomrule
\end{tabular}
\caption{Technologie-Stack}
\end{reporttable}

% ================================================================== 4 Kern
\section{Kern-Interpreter}\label{sec:kern}
Dieses Kapitel beschreibt die Phasen des Kern-Interpreters in der Reihenfolge, in der ein Programm sie durchläuft.
Der Schwerpunkt liegt auf den Entwurfsentscheidungen, die nicht unmittelbar aus dem Code ersichtlich sind.

\subsection{Lexikalische Analyse}
Der Lexer wird aus dem Lexer-Teil von \code{interpreter/src/main/antlr4/MiniCpp.g4} generiert. Schlüsselwörter
(einschließlich \code{new}, \code{delete} und \code{nullptr}) sind eigene Token-Typen und stehen vor der
\code{Identifier}-Regel, sodass ANTLR bei gleicher Länge das Schlüsselwort bevorzugt. Leerraum, Zeilen- und
Blockkommentare sowie Präprozessorzeilen (\code{\#include <iostream>}) werden verworfen -- so lassen sich
MiniC++-Programme unverändert auch mit einem C++-Compiler übersetzen. Zeichen- und Zeichenkettenliterale erlauben die
Escape-Sequenzen \code{\bs b \bs t \bs n \bs r \bs\dq{} \bs\sq{} \bs\bs{} \bs 0}. Jedes Token trägt Zeile und Spalte;
diese Positionen wandern über den AST bis in die Diagnosen.

\subsection{Syntaxanalyse}
Der Parser besitzt zwei Startregeln: \code{program} für Dateien (nur Deklarationen) und \code{replInput} für die
REPL, die Deklarationen und Anweisungen in beliebiger Reihenfolge und einen abschließenden Ausdruck ohne Semikolon
erlaubt. Die Ausdrucksgrammatik ist als Präzedenzkaskade formuliert -- eine Regel pro Präzedenzstufe --, was die
Präzedenz explizit und die Grammatik frei von Mehrdeutigkeiten hält:

\begin{lstlisting}[caption={Auszug aus der Parser-Grammatik}]
assignmentExpr     : logicalOrExpr (ASSIGN assignmentExpr)? ;        // rechtsassoziativ
logicalOrExpr      : logicalAndExpr (OROR logicalAndExpr)* ;
  ...
multiplicativeExpr : unaryExpr ((STAR | SLASH | PERCENT) unaryExpr)* ;
unaryExpr          : (NOT | PLUS | MINUS | STAR | AMP) unaryExpr | newExpr | postfixExpr ;
postfixExpr        : primaryExpr postfixPart* ;
postfixPart        : LPAREN argList? RPAREN | DOT Identifier | ARROW Identifier ;

type               : (baseType | Identifier) STAR* ;                 // 'T', 'T*', 'T**'
statement          : block | varDecl SEMI | ifStmt | whileStmt
                   | returnStmt SEMI | deleteStmt SEMI | expr SEMI ;
\end{lstlisting}

Da \code{varDecl} in \code{statement} vor \code{expr} steht und ANTLR bei Mehrdeutigkeit die erste passende
Alternative wählt, wird \code{a * b;} -- wie in C++ -- als Zeigerdeklaration erkannt. Die Grammatik akzeptiert bewusst
etwas mehr als die Sprache (z.\,B. Klassenbezeichner an beliebiger Typposition); die Einschränkungen prüfen Resolver
und Typprüfung, die präzisere Fehlermeldungen liefern können als der Parser.

\subsubsection*{Fehlerbehandlung}
Die Standard-Fehlerausgabe von ANTLR wird durch einen \code{SyntaxErrorCollector} ersetzt, der jeden Fehler als
\code{Diagnostic} mit der Länge des betroffenen Tokens aufzeichnet. Zusätzlich merkt er sich, ob ein Fehler \emph{am
Dateiende} auftrat. Dieses Flag (\code{ParseResult.incomplete}) nutzt die REPL, um unvollständige Eingaben -- etwa
eine offene geschweifte Klammer -- von echten Syntaxfehlern zu unterscheiden. Bei Syntaxfehlern wird kein AST
erzeugt; die Fehlerbehandlung von ANTLR sorgt aber dafür, dass alle Syntaxfehler einer Datei in einem Durchlauf
gemeldet werden.

\subsection{Abstrakter Syntaxbaum}
Der \code{ASTBuildVisitor} übersetzt den Parse-Tree in einen handgeschriebenen AST. Dabei werden syntaktische Details
entfernt (Klammern, Präzedenzstufen), Kurzformen normalisiert (\code{T x(a)} wird zu einer Deklaration mit
Konstruktoraufruf) und jedem Knoten sein vollständiger Quellbereich mitgegeben. Die Knoten gliedern sich in vier
Hierarchien:

\begin{reporttable}
\begin{tabular}{L{0.16}L{0.76}}
\thead \hcell{Basisklasse} & \hcell{Knoten} \\ \midrule
\code{Decl} & \code{ClassDecl}, \code{ConstructorDecl}, \code{FieldDecl}, \code{FunctionDecl}, \code{MethodDecl},
  \code{ParameterDecl}, \code{VariableDecl} \\
\code{Stmt} & \code{BlockStmt}, \code{VariableStmt}, \code{ExprStmt}, \code{IfStmt}, \code{WhileStmt},
  \code{ReturnStmt}, \code{DeleteStmt} \\
\code{Expr} & \code{AssignExpr}, \code{BinaryExpr}, \code{UnaryExpr}, \code{CallExpr}, \code{MemberAccessExpr},
  \code{NewExpr}, \code{VarExpr}, Literale (\code{Int}, \code{Bool}, \code{Char}, \code{String}, \code{Nullptr}),
  \code{ErrorExpr} \\
\code{Type} & \code{PrimitiveType}, \code{ClassType}, \code{PointerType}, \code{ReferenceType}, \code{NullptrType} \\
\bottomrule
\end{tabular}
\caption{Knotenhierarchien des AST}
\end{reporttable}

Wurzelknoten sind \code{Program} (Datei) und \code{ReplInput} (REPL-Eingabe). Spätere Phasen annotieren den AST,
statt eigene Tabellen aufzubauen: Der Resolver setzt an \code{VarExpr} die referenzierte Deklaration, die Typprüfung
setzt an jedem \code{Expr} den inferierten Typ und an \code{CallExpr} bzw.\ \code{NewExpr} die gewählte Überladung.
Genau diese Annotationen nutzt der Language-Server später für Hover, Go-to-Definition und Umbenennen.

\subsection{Namensauflösung}
Der \code{ASTResolveVisitor} arbeitet in zwei Durchläufen. \textbf{Pass 1} sammelt alle Klassen und Funktionen im
\code{GlobalScope}, verknüpft jede Klasse mit ihrer Basisklasse, erkennt zyklische Vererbung und prüft
Member-Deklarationen (doppelte Felder, erneut deklarierte Felder in abgeleiteten Klassen, doppelte Signaturen). Weil
alle globalen Namen danach bekannt sind, dürfen Funktionen und Klassen vor ihrer Definition verwendet werden.
\textbf{Pass 2} durchläuft alle Rümpfe mit einem Stapel lokaler Gültigkeitsbereiche und bindet jeden Bezeichner in
der Reihenfolge \emph{lokal $\rightarrow$ eigene Member $\rightarrow$ geerbte Member $\rightarrow$ global}. Da lokale
Namen erst beim Erreichen ihrer Deklaration in den aktuellen Bereich eingetragen werden, entsteht für Variablen
automatisch die define-before-use-Semantik.

\begin{lstlisting}[language=C++,caption={Typische Fehler der Namensauflösung (aus \code{ErrorTest})}]
int main() { x = 1; int x; }            // 1:14: error: use of undeclared identifier 'x'
int main() { int x; bool x; }           // error: redeclaration of 'x'
class A : public B { public: };         // error: unknown base class 'B'
class A : public B { public: }; class B : public A { public: };   // error: 'A' inherits from itself
void f(int a) { } void f(int b) { }     // error: redefinition of function 'f(int)'
\end{lstlisting}

Für die REPL stellt der Resolver zusätzlich Einzeloperationen bereit (\code{declareClass}, \code{declareFunction},
\code{resolveSessionStatement}), mit denen Eingaben einzeln und damit stets define-before-use aufgelöst werden.

\subsection{Typprüfung}
Da \code{auto} nicht zum Sprachumfang gehört, ist keine Typinferenz im Sinne von Hindley-Milner nötig; jede Variable
trägt ihren deklarierten Typ. Der \code{TypeCheckVisitor} bestimmt die Typen von Ausdrücken bottom-up und prüft dabei
die Regeln aus Kapitel~\ref{sec:sprache}: gleiche Operandentypen, zulässige Operatoren pro Typ, L-Werte links von
\code{=} und bei Referenzparametern, Bedingungstypen, Zeigerkompatibilität und Konstruktoraufrufe. Darüber hinaus
prüft er:
\begin{itemize}
  \item \textbf{Return-Pfade:} Eine Nicht-\code{void}-Funktion muss auf allen Pfaden ein \code{return} erreichen
    (\code{alwaysReturns} analysiert \code{if}/\code{else}-Verzweigungen und Blöcke; \code{int main()} ist
    ausgenommen).
  \item \textbf{Overrides:} Eine Methode mit gleicher Signatur wie eine virtuelle Basismethode muss denselben
    Rückgabetyp haben und wird selbst als virtuell markiert (\code{isEffectivelyVirtual}).
  \item \textbf{Rekursive Einbettung:} Eine Klasse darf sich nicht selbst als Wertfeld enthalten, auch nicht indirekt
    (Zeigerfelder sind erlaubt).
  \item \textbf{Einstiegspunkt:} Es muss genau ein parameterloses \code{main} mit Rückgabetyp \code{int} oder
    \code{void} existieren.
\end{itemize}

\subsubsection*{Überladungsauflösung}
Die Auflösung von Funktions-, Methoden- und Konstruktoraufrufen erfolgt zentral in \code{pickOverload}. Für jeden
Kandidaten berechnet \code{matchCost} die Anzahl der nötigen impliziten Konvertierungen ($-1$ = nicht aufrufbar). Ein
Referenzparameter verlangt ein L-Wert-Argument und darf nicht an einen konvertierten Zeiger binden, da dieser ein
temporärer Wert wäre.

\begin{lstlisting}[language=Java,caption={Kern der Überladungsauflösung in \code{TypeCheckVisitor}}]
List<D> exact = ..., converting = ...;
for (D candidate : candidates) {
  int cost = matchCost(params.apply(candidate), args, argTypes);
  if (cost == 0) exact.add(candidate);
  else if (cost > 0) converting.add(candidate);
}
List<D> best = !exact.isEmpty() ? exact : converting;
if (best.size() == 1) return best.get(0);
// sonst: "no matching function for call to 'f(int, bool)'; candidates are: ..."
//    bzw. "call to 'f(B*)' is ambiguous; candidates are: ..."
\end{lstlisting}

\subsection{Laufzeitmodell}\label{sec:laufzeit}
Der \code{Interpreter} ist ein Tree-Walking-Interpreter: Er implementiert \code{AstVisitor<Object>}, wobei Ausdrücke
zu Laufzeitwerten und Anweisungen zu \code{null} oder einem \code{Returned}-Signal auswerten. Die Rückgabe über ein
Signalobjekt statt über eine Java-Exception hält den Aufrufpfad frei von teurem Stack-Unwinding. Laufzeitwerte sind
\code{Integer}, \code{Boolean}, \code{Character}, \code{String}, \code{Pointer} oder \code{ObjectValue}.

\begin{figure}[htbp]
\centering
\begin{tikzpicture}[font=\footnotesize]
  % Frame
  \draw[navy, fill=white, rounded corners=2pt, line width=0.6pt] (0,0) rectangle (4.4,5.8);
  \node[font=\footnotesize\bfseries, color=navy] at (2.2,5.45) {Frame (Activation Record)};
  \foreach \t/\y/\n in {int x/4.75/fx, Konto a/3.85/fa, int\& r\ \ (Alias)/2.95/fr, Konto* p/2.05/fp}
    \node[dbox, fill=zebra, draw=rule, minimum width=3.8cm, minimum height=0.6cm] (\n) at (2.2,\y) {\t};
  \node[lbl] at (2.2,0.45) {\code{Map<Decl, Cell>} + owned-Liste};
  % Mitte
  \node[dbox, minimum width=2.6cm, minimum height=0.7cm] (c42) at (7.3,5.0) {Cell: 42};
  \node[lbl] at (7.3,4.35) {r = Alias dieser Cell};
  \draw[navy, fill=headbg, rounded corners=2pt, line width=0.6pt] (5.6,1.9) rectangle (10.4,3.95);
  \node[font=\footnotesize\bfseries] at (8.0,3.6) {Cell $\rightarrow$ ObjectValue (Konto)};
  \node[dbox, fill=white, draw=rule, minimum width=2.0cm, minimum height=0.95cm] at (6.85,2.6) {owner\\[-1pt]{\scriptsize\itshape\color{ink2}Cell}};
  \node[dbox, fill=white, draw=rule, minimum width=2.0cm, minimum height=0.95cm] at (9.15,2.6) {balance\\[-1pt]{\scriptsize\itshape\color{ink2}Cell}};
  \coordinate (obj) at (5.6,3.0);
  \node[dbox, fill=white, draw=accent, minimum width=3.6cm, minimum height=0.95cm] (ptr) at (7.4,0.65)
    {Pointer(target)\\[-1pt]{\scriptsize\itshape\color{ink2}record; NULL = nullptr}};
  \draw[arr] (fx.east) -- (c42.west);
  \draw[arr] (fa.east) -- (obj);
  \draw[arr, dashed] (fr.east) -- (c42.south west);
  \draw[arr] (fp.east) -- (ptr.west);
  % rechts: vtable und Heap
  \draw[ink2, fill=white, rounded corners=2pt, line width=0.6pt] (11.4,3.6) rectangle (16,5.8);
  \node[font=\footnotesize\bfseries, color=ink2] at (13.7,5.45) {vtable(Konto)};
  \node[anchor=west, font=\ttfamily\scriptsize] at (11.55,4.9) {gebuehr()\ \ \ \ \ \ $\rightarrow$ Konto};
  \node[anchor=west, font=\ttfamily\scriptsize] at (11.55,4.5) {einzahlen(int) $\rightarrow$ Konto};
  \node[lbl, anchor=west] at (11.55,3.95) {\code{Map<String, MethodDecl>}};
  \draw[ink2, fill=white, rounded corners=2pt, line width=0.6pt] (11.4,0) rectangle (16,2.9);
  \node[font=\footnotesize\bfseries, color=ink2] at (13.7,2.55) {Heap (new)};
  \node[dbox, minimum width=4.0cm, minimum height=1.3cm] (hc) at (13.7,1.1)
    {Cell (heap = true)\\[-1pt]{\scriptsize\itshape\color{ink2}$\rightarrow$ ObjectValue}};
  \draw[arr, accent] (ptr.east) -- (hc.west);
  \draw[arr, dashed] (10.4,3.6) -- (11.4,4.4) node[lbl, midway, above left] {dynamicClass};
\end{tikzpicture}
\caption{Speichermodell: Frames, Cells, Objekte, Zeiger und vtables}
\end{figure}

\subsubsection*{Cells als Speicherorte}
Zentrale Abstraktion ist die \code{Cell}, ein veränderlicher Speicherort. Jede Variable, jeder By-Value-Parameter und
jedes Feld besitzt genau eine Cell. Eine \textbf{Referenz} ist schlicht ein zweiter Name für dieselbe Cell --
Zuweisungen über die Referenz ändern damit automatisch das Original. Ein \textbf{Zeiger} ist der Record
\code{Pointer(Cell target)}; \code{nullptr} ist \code{Pointer(null)}, und zwei Zeiger sind gleich, wenn sie auf
dieselbe Cell zeigen.

\subsubsection*{Lebensdauer und Fehlererkennung}
Jeder Frame (Activation Record) führt eine Liste der Cells, die ihm gehören. Endet ein Block oder eine Funktion,
werden diese Cells per \code{kill()} als tot markiert -- rekursiv einschließlich der Felder enthaltener Objekte.
Heap-Cells aus \code{new} tragen zusätzlich das Flag \code{heap}; \code{delete} prüft, dass der Zeiger auf eine
lebende Heap-Cell zeigt, und tötet sie. Jeder Zugriff über einen Zeiger oder eine Referenz prüft den Zustand der
Ziel-Cell. So werden \code{nullptr}-Dereferenzierung, hängende Zeiger, Use-after-free, doppeltes \code{delete} und
\code{delete} auf Stack-Objekte zuverlässig und mit Quellposition erkannt, ohne eine Speicherverwaltung zu simulieren
-- die eigentliche Freigabe übernimmt der Garbage Collector der JVM.

\subsubsection*{Objekte, Kopien und Slicing}
Ein \code{ObjectValue} hält seine Felder in einer \code{LinkedHashMap<String, Cell>} in Deklarationsreihenfolge.
Zuweisungen und Wertparameter kopieren Objekte tief (\code{copy}); ist der statische Zieltyp eine Basisklasse,
entsteht mit \code{sliceTo} eine Kopie, die nur deren Felder enthält. Temporäre Objekte (etwa Rückgabewerte) werden
nicht erneut kopiert.

\subsubsection*{Virtueller Dispatch}
Für jede Klasse wird bei Bedarf eine vtable als \code{Map<String, MethodDecl>} aufgebaut: Sie übernimmt die Einträge
der Basisklasse und überschreibt sie mit den eigenen Methoden. Schlüssel ist der Methodenname zusammen mit den
Parametertypen, sodass überladene virtuelle Methoden getrennt bleiben. Nicht-virtuelle Aufrufe werden statisch
gebunden, virtuelle über die vtable der \emph{dynamischen Klasse} des Empfängers. Diese ist normalerweise die Klasse
des Objekts, wird aber während eines Basisklassen-Konstruktors vorübergehend auf die Basisklasse gesetzt -- das bildet
die C++-Regel nach, dass im Konstruktor noch nicht in abgeleitete Klassen dispatcht wird.

\begin{lstlisting}[language=Java,caption={Dynamischer Dispatch im Interpreter}]
private MethodDecl dispatch(MethodDecl method, ObjectValue receiver) {
  if (!method.isEffectivelyVirtual()) return method;
  MethodDecl override = vtable(receiver.getDynamicClass()).get(vtableKey(method));
  return override != null ? override : method;
}
\end{lstlisting}

\subsubsection*{Rekursionstiefe}
Ein Tree-Walker benötigt pro interpretiertem Aufruf mehrere Java-Stackframes. Damit auch tiefe Rekursion (z.\,B.\
100\,000 Ebenen) funktioniert, führt \code{MiniCpp.onLargeStack} den Interpreter auf einem eigenen Thread mit 512 MiB
Stack aus. Ein Zähler begrenzt die Aufruftiefe auf \code{MAX\_CALL\_DEPTH = 100\_000}; verbleibende
\code{StackOverflowError}s (etwa durch extrem tief verschachtelte Ausdrücke) werden in einen regulären Laufzeitfehler
übersetzt.

\subsection{REPL}
Die REPL (\code{repl.Repl}) lädt optional eine Datei, führt deren \code{main} im \emph{Sitzungs-Scope} aus und hält
diesen offen, sodass die lokalen Variablen von \code{main} anschließend zur Verfügung stehen. Jede Eingabe wird mit
der Startregel \code{replInput} geparst. Meldet der Parser einen Fehler am Eingabeende, zeigt die REPL den
Fortsetzungsprompt \code{...>} und sammelt weitere Zeilen. Eine vollständige Eingabe wird als Ganzes aufgelöst,
geprüft und ausgeführt; vor der Verarbeitung wird mit \code{GlobalScope.snapshot()} ein Sicherungspunkt angelegt, auf
den bei einem Fehler zurückgesetzt wird. Eine fehlerhafte Eingabe hinterlässt damit keine halb definierten Klassen
oder Funktionen.

\begin{lstlisting}[caption={REPL-Sitzung}]
minicpp> int x = 6;
minicpp> int sq(int n) {
     ...>   return n * n;
     ...> }
minicpp> sq(x) + 1
37
minicpp> :vars
int x = 6
\end{lstlisting}

Neue Variablen landen im Sitzungs-Scope, neue Funktionen und Klassen im globalen Scope. Da MiniC++ keine globalen
Variablen kennt, sind Sitzungsvariablen in Funktionen nicht sichtbar. Befehle: \code{:vars}, \code{:functions},
\code{:classes}, \code{:load <datei>}, \code{:reset}, \code{:cancel}, \code{:help}, \code{:quit}.

\subsection{Kommandozeile und C++-Export}\label{sec:cppexport}
Die Klasse \code{Main} stellt die Unterbefehle \code{run}, \code{check}, \code{repl}, \code{ast} und \code{to-cpp}
bereit (Anhang~\ref{app:build}). Der Exit-Code von \code{run} ist der Rückgabewert von \code{main}, bei Fehlern 1.

Der \code{CppExporter} übersetzt ein geprüftes MiniC++-Programm in Standard-C++17 mit identischem Verhalten. Er ist
die Grundlage des differenziellen Tests gegen GCC (Abschnitt~\ref{sec:gcc}) und überbrückt die Stellen, an denen
MiniC++ bewusst von C++ abweicht:
\begin{itemize}
  \item define-after-use: Klassen werden vorwärts deklariert und nach Abhängigkeiten sortiert, Funktionen erhalten
    Prototypen, Methoden werden nach allen Klassen außerhalb definiert;
  \item Standardwerte für nicht initialisierte Variablen und Felder, \code{new T} wird zu \code{new T()};
  \item Zeichenkettenliterale erhalten den Typ \code{string}, die \code{print\_*}-Funktionen werden in einem Präludium
    definiert, \code{void main()} wird gekapselt;
  \item Basisklassen erhalten einen virtuellen Destruktor, damit \code{delete} über einen Basiszeiger definiert ist.
\end{itemize}

% ================================================================== 5 LSP
\section{Language-Server}\label{sec:lsp}
Der Language-Server macht MiniC++ in jedem LSP-fähigen Editor nutzbar; Referenz-Client ist eine VS-Code-Extension. Er
ist in Java mit Eclipse LSP4J implementiert, kommuniziert per JSON-RPC über stdin/stdout und verwendet für jede
Analyse exakt dieselbe Pipeline wie der Interpreter. Damit sind Diagnosen im Editor und auf der Kommandozeile
garantiert identisch.

\subsection{Architektur}
\begin{reporttable}
\begin{tabular}{L{0.4}L{0.52}}
\thead \hcell{Klasse} & \hcell{Verantwortung} \\ \midrule
\code{MiniCppLanguageServerMain} & Start über stdio, Verbindung mit dem Client \\
\code{MiniCppLanguageServer} & Lebenszyklus (\code{initialize}, \code{shutdown}, \code{exit}), Capabilities \\
\code{MiniCppTextDocumentService} & Dokumentsynchronisation, entprellte Diagnosen, Verteilung der Anfragen \\
\code{Document}, \code{Analysis} & offenes Dokument; unveränderlicher Analyse-Snapshot einer Version \\
\code{SymbolIndex} & jedes Vorkommen eines Bezeichners mit der Deklaration, auf die es verweist \\
\code{AstNodes}, \code{Names} & AST-Traversierung, Gültigkeitsbereich an einer Position, Signaturen \\
\code{Hovers}, \code{Completions}, \code{Navigation}, \code{CodeFormatter}, \code{CodeActions} & die einzelnen
  Features \\
\code{SourceText}, \code{Tokens}, \code{Positions} & Offsets $\leftrightarrow$ Positionen (Interpreter 1-basiert,
  LSP 0-basiert), Token-Suche, Diagnose-Umrechnung \\ \bottomrule
\end{tabular}
\caption{Klassen des Language-Servers}
\end{reporttable}

Beim \code{initialize} meldet der Server folgende Fähigkeiten: inkrementelle Textsynchronisation, Hover, Completion
mit den Triggerzeichen \code{.} und \code{>}, Definition, Referenzen, Document Highlight, Document Symbol, Rename mit
\code{prepareRename}, Formatierung und Code-Actions der Art \code{quickfix}.

\subsection{Synchronisation und Entprellung}\label{sec:sync}
Der Client überträgt Änderungen inkrementell als Bereichsänderungen, die \code{Document.update} der Reihe nach auf den
Text anwendet. Jede Änderung plant auf einem einzelnen Analyse-Thread eine Neuanalyse mit 200\,ms Verzögerung und
bricht eine noch ausstehende ab. Beim Tippen wird so nur einmal nach der letzten Änderung analysiert. Ergebnisse
werden nur veröffentlicht, wenn die analysierte Version noch aktuell ist.

Anfragen wie Hover oder Completion warten nicht auf den Timer: \code{Document.analysis()} berechnet die Analyse
sofort, falls für die aktuelle Version noch keine vorliegt. Eine Anfrage sieht dadurch nie einen veralteten Stand.

\begin{figure}[htbp]
\centering
\begin{tikzpicture}[msg/.style={font=\scriptsize\itshape, color=ink2, fill=white, inner sep=1pt}]
  \foreach \t/\x [count=\i from 0] in
     {VS Code\\(vscode-languageclient)/2.0, MiniCppTextDocument-\\Service/6.05,
      Analyzer-Thread\\(ScheduledExecutor)/10.1, MiniCpp.compile\\+ SymbolIndex/14.15}
  {
    \node[dbox, text width=3.6cm, minimum height=0.95cm, font=\scriptsize\bfseries] (l\i) at (\x,0) {\t};
    \draw[rule, dashed, line width=0.7pt] (l\i.south) -- (\x,-6.4);
  }
  \newcommand{\msg}[5][]{%
    \draw[arr, #1] ($(#3,#2)+(0.05,0)$) -- ($(#4,#2)-(0.05,0)$)
      node[msg, midway, above=1pt] {#5};}
  \msg{-1.1}{2.0}{6.05}{didChange (Bereichsänderung, v7)}
  \msg{-1.9}{6.05}{10.1}{schedule(200 ms), vorige Aufgabe cancel}
  \msg{-2.7}{2.0}{6.05}{didChange (v8) $\rightarrow$ erneut entprellen}
  \msg{-3.6}{10.1}{14.15}{Document.analysis()\ \ (v8)}
  \draw[arr, dashed] (14.1,-4.4) -- (10.15,-4.4) node[msg, midway, above=1pt] {Analysis(Tokens, Compilation, Index)};
  \draw[arr, accent] (10.05,-5.2) -- (2.05,-5.2) node[msg, midway, above=1pt, text=accent] {publishDiagnostics (Version 8)};
  \msg{-6.0}{2.0}{6.05}{hover / completion / \dots{} $\rightarrow$ Analyse sofort, falls ausstehend}
\end{tikzpicture}
\caption{Ablauf von Änderungen, Entprellung und Anfragen}
\end{figure}

\begin{hinweis}
Die Synchronisation ist inkrementell, die Analyse selbst aber nicht: Nach jeder (entprellten) Änderung läuft die
vollständige Pipeline über das gesamte Dokument. Bei den gemessenen Frontend-Zeiten (Abschnitt~\ref{sec:perf}:
ca.\ 3\,ms für 100 Klassen) ist das für realistische Dateigrößen unkritisch. Im Arbeitsplan war ursprünglich
inkrementelles \emph{Parsen} vorgesehen -- im Bericht sollte die Abweichung begründet werden.
\end{hinweis}

\subsection{Analyse-Snapshot und Fehlertoleranz}
Eine \code{Analysis} bündelt für eine Dokumentversion den Quelltext, den Token-Strom, die \code{Compilation} und den
Symbolindex. Zwei Mechanismen machen den Server robust gegenüber unfertigem Code:
\begin{itemize}
  \item \textbf{Letzte parsebare Version.} Während des Tippens ist Code oft syntaktisch unvollständig -- nach
    \code{konto.} existiert kein AST. Jede Analyse verweist daher über \code{lastParsed} auf die jüngste Version mit
    AST. Die Completion ermittelt Typ und Member des Ausdrucks vor dem Punkt aus dieser Version, und zwar auch über
    Ketten wie \code{a.b->c().}.
  \item \textbf{Absturzisolation.} Ausnahmen in Pipeline oder Indexierung werden abgefangen und protokolliert; der
    Server liefert dann leere Ergebnisse, statt die Verbindung zu verlieren.
\end{itemize}

\subsection{Symbolindex}
AST-Knoten kennen ihren Quellbereich, aber nicht die exakte Position ihres \emph{Namens}. Der \code{SymbolIndex}
kombiniert deshalb den aufgelösten AST mit dem Token-Strom und legt für jedes Vorkommen eines Bezeichners einen
Eintrag \code{Occurrence(start, end, symbol, target, declaration)} in einer nach Offset sortierten \code{TreeMap} ab.
Eine Positionsanfrage ist damit ein \code{floorEntry}-Aufruf in $O(\log n)$.

Die Unterscheidung zwischen \code{target} und \code{symbol} ist der Kern der Navigationsfeatures: \code{target} ist
die exakte Deklaration (z.\,B. die gewählte Überladung oder der aufgerufene Konstruktor) und dient Go-to-Definition
und Hover. \code{symbol} gruppiert alle Vorkommen, die gemeinsam umbenannt werden müssen: Konstruktoren gehören zu
ihrer Klasse, überschreibende Methoden zur überschriebenen Basismethode.

\subsection{Features}
\begin{reporttable}
\begin{tabular}{L{0.28}L{0.64}}
\thead \hcell{Feature} & \hcell{Verhalten} \\ \midrule
\code{publishDiagnostics} & Fehler aller statischen Phasen (Syntax, Namen, Typen) mit exaktem Bereich \\
\code{hover} & Deklaration als C++-Signatur mit Art (\code{local variable}, \code{virtual method},
  \code{overrides B::f} \dots); an Operatoren und Literalen der statische Typ des Ausdrucks \\
\code{completion} & Variablen und Parameter im Gültigkeitsbereich, Member der eigenen Klasse, Funktionen, Klassen,
  Schlüsselwörter, Typen; nach \code{.}/\code{->} die Member (inkl.\ geerbter) des Objekts \\
\code{definition} & gewählte Überladung, statisch gebundene Methode, Konstruktor bei \code{A(\dots)} /
  \code{new A} \\
\code{references}, \code{documentHighlight}, \code{rename} & Klassen samt Konstruktornamen und Typverwendungen;
  überschreibende Methoden gemeinsam mit der Basismethode; Built-ins und Keywords sind nicht umbenennbar \\
\code{documentSymbol} & Outline mit Klassen, Feldern, Konstruktoren, Methoden und Funktionen \\
\code{formatting} & auf Basis des Parse-Trees: K\&R-Klammern, \code{public:} auf Klassenebene, \code{T* p},
  Leerzeichen um binäre Operatoren; Kommentare und einzelne Leerzeilen bleiben erhalten \\
\code{codeAction} & Quick Fixes: fehlendes Token einfügen, überzähliges entfernen, ähnlichen Namen vorschlagen,
  \code{.} $\rightarrow$ \code{->}, \code{()} an Methodennamen, leeres \code{main} ergänzen \\ \bottomrule
\end{tabular}
\caption{Umgesetzte LSP-Features}
\end{reporttable}

\begin{todo}
Screenshots der Features in VS Code (Hover, Completion nach \code{->}, Quick Fix, Rename) einfügen.
\end{todo}

\subsection{VS-Code-Extension}
Die Extension (\code{vscode/}) registriert die Sprache \code{minicpp} für \code{*.mcpp}-Dateien, liefert eine
TextMate-Grammatik für Syntaxhervorhebung sowie eine Sprachkonfiguration (Kommentare, Klammerpaare) und startet den
Server über \code{vscode-languageclient}. Der Server wird in folgender Reihenfolge gesucht: konfigurierter Pfad
(\code{minicpp.server.path}), im \code{.vsix}-Paket gebündelter Server (direkt mit \code{java} gestartet, da
Startskripte im Paket ihr Ausführungsrecht verlieren), zuletzt der im Repository gebaute Server. Das Java-Programm
lässt sich mit \code{minicpp.java.path} festlegen, \code{minicpp.trace.server} protokolliert den JSON-RPC-Verkehr.

% ================================================================== 6 GenAI
\section{GenAI-Assistenzserver}\label{sec:genai}
Der Assistenzserver (Modul \code{mcp}) stellt KI-gestützte Code-Assistenz für MiniC++ als HTTP/JSON-Dienst bereit. Er
nutzt ein lokal laufendes Sprachmodell über Ollama und verbindet dessen Antworten mit dem echten MiniC++-Compiler.

\begin{todo}
Das Modul heißt \code{mcp}, implementiert aber eine REST-Schnittstelle und nicht das \emph{Model Context Protocol}
(JSON-RPC mit Tools/Resources). Für die Endfassung klären: entweder Begriff im Bericht präzisieren
(\enquote{GenAI-REST-Server}) oder eine MCP-Schnittstelle ergänzen.
\end{todo}

\subsection{Endpunkte}
\begin{reporttable}
\begin{tabular}{L{0.19}L{0.36}L{0.37}}
\thead \hcell{Endpunkt} & \hcell{Anfrage (JSON)} & \hcell{Antwort} \\ \midrule
\code{POST /complete} & \code{code}, \code{line} (1-basiert), \code{column} (0-basiert) & \code{completion}:
  einzufügender Text \\
\code{POST /explain} & \code{code} & \code{explanation} \\
\code{POST /refactor} & \code{code}, optional \code{instruction} & \code{code}, \code{rationale}, \code{compiles},
  \code{diagnostics} \\
\code{POST /detect-bugs} & \code{code} & \code{compilerDiagnostics}, \code{findings[\{line, description\}]},
  \code{analysis} \\
\code{GET /health} & -- & Status, Provider, Modell, Erreichbarkeit \\ \bottomrule
\end{tabular}
\caption{Endpunkte des Assistenzservers}
\end{reporttable}

Der Server basiert auf dem im JDK enthaltenen \code{com.sun.net.httpserver.HttpServer} mit einem Pool von vier
Threads und Gson für JSON; weitere Abhängigkeiten gibt es nicht. Er bindet standardmäßig nur an
\code{127.0.0.1:8080}. Fehler werden einheitlich als \code{\{\dq error\dq: \dots\}} gemeldet: HTTP 400 bei
ungültiger Anfrage, 404 bei unbekanntem Pfad, 405 bei falscher Methode, 413 bei Anfragen über 1\,MiB und 502, wenn
das Modell nicht erreichbar ist.

\subsection{Provider-Abstraktion}
Die Schnittstelle \code{GenAiProvider} kapselt das Sprachmodell (\code{generate(system, prompt)},
\code{isAvailable()}). Der \code{OllamaProvider} ruft \code{/api/generate} ohne Streaming und mit Temperatur 0,2 auf
-- Code-Assistenz soll fokussiert statt kreativ sein. Modell (Standard \code{codellama}), URL und Timeout (120\,s)
sind per Kommandozeile oder Umgebungsvariable konfigurierbar. Der \code{MockProvider} implementiert dieselbe
Schnittstelle deterministisch; er ermöglicht Tests des gesamten Servers ohne laufendes Modell und einen Betrieb mit
\code{--provider mock}.

\subsection{Prompt-Design}
Sprachmodelle kennen C++, aber nicht MiniC++. Jeder System-Prompt enthält deshalb eine kompakte Sprachbeschreibung,
die neben den Features ausdrücklich aufzählt, was fehlt (Arrays, Casts, \code{++}, \code{+=}, \code{this},
Standardbibliothek \dots). Ohne diese Negativliste schlagen Modelle bevorzugt idiomatisches C++ vor, das MiniC++
nicht akzeptiert. Die endpunktspezifischen Anweisungen erzwingen ein maschinenlesbares Antwortformat:
\begin{itemize}
  \item \code{/complete} markiert die Cursorposition mit \code{<CURSOR>} und verlangt ausschließlich den
    einzufügenden Text; umschließende Markdown-Codeblöcke werden trotzdem defensiv entfernt.
  \item \code{/refactor} verlangt den vollständigen Code in genau einem \code{\textasciigrave\textasciigrave\textasciigrave cpp}-Block, gefolgt von
    einer Zeile \code{Rationale:}.
  \item \code{/detect-bugs} übergibt den Code mit Zeilennummern und verlangt Befunde im Format
    \code{LINE <n>: <text>} oder \code{NONE}; ein regulärer Ausdruck überführt sie in strukturierte \code{findings}.
\end{itemize}

\subsection{Absicherung durch den Compiler}
Das wichtigste Qualitätsmerkmal des Servers ist, dass er Modellantworten nicht ungeprüft weitergibt:
\begin{itemize}
  \item \textbf{Grounding:} Für \code{/explain} und \code{/detect-bugs} werden die Diagnosen des MiniC++-Compilers an
    den Prompt angehängt. Das Modell muss Syntax- und Typfehler nicht raten, sondern kann sie als gesichert
    behandeln.
  \item \textbf{Verifikation:} Jeder Refactoring-Vorschlag wird vor der Rückgabe kompiliert; das Feld
    \code{compiles} und die Diagnosen zeigen dem Aufrufer, ob der Vorschlag gültiges MiniC++ ist.
  \item \textbf{Strukturierte Ergänzung:} \code{/detect-bugs} liefert die Compilerdiagnosen zusätzlich getrennt von
    den Modellbefunden, sodass ein Client sichere und heuristische Befunde unterschiedlich darstellen kann.
\end{itemize}
Code ohne \code{main} wird dabei nur als Deklarationssammlung geprüft, damit auch einzelne Funktionen oder Klassen
analysiert werden können.

\begin{todo}
Qualitative Evaluation der Modellantworten ergänzen (z.\,B. Anteil kompilierender Refactorings für \code{codellama}
vs.\ \code{deepseek-coder} auf den Testprogrammen, typische Fehlermuster, Antwortzeiten).
\end{todo}

% ================================================================== 7 Evaluation
\section{Qualitätssicherung und Evaluation}\label{sec:evaluation}
\subsection{Teststrategie}
Alle 218 automatisierten Tests laufen mit \code{./gradlew build} bzw.\ \code{./gradlew test} und waren zum
Redaktionsschluss grün. Sie gliedern sich wie folgt:

\begin{reporttable}
\begin{tabular}{L{0.24}L{0.12}>{\raggedleft\arraybackslash}p{0.06\linewidth}L{0.48}}
\thead \hcell{Testklasse} & \hcell{Modul} & \hcell{Tests} & \hcell{Gegenstand} \\ \midrule
\code{ErrorTest} & interpreter & 129 & ungültige Programme: erwartete Phase und Meldung des ersten Fehlers, inkl.\
  Laufzeitfehler \\
\code{CppExporterTest} & interpreter & 14 & Übersetzung nach C++ \\
\code{ReplTest} & interpreter & 13 & Sitzungen, Fortsetzungszeilen, Rücksetzen bei Fehlern, Befehle \\
\code{ParserTest} & interpreter & 10 & Präzedenz, Literale und Escapes, Kurzformen, Referenztypen, Klassen \\
\code{ProgramTest} & interpreter & 9 & Golden-File-Tests der Beispielprogramme \\
\code{FeaturesTest} & lsp & 19 & Hover, Completion, Navigation, Rename, Formatierung, Quick Fixes \\
\code{DiagnosticsTest} & lsp & 3 & Diagnosen Ende-zu-Ende mit Client-Stub: Positionen, Syntaxfehler, Löschen beim
  Schließen \\
\code{AssistantServiceTest} & mcp & 9 & Prompts, Antwortverarbeitung, Compiler-Absicherung (Mock) \\
\code{McpServerTest} & mcp & 8 & HTTP-Ebene, Fehlercodes \\
\code{OllamaProviderTest} & mcp & 4 & Ollama-Protokoll gegen lokalen Test-Server \\ \midrule
\textbf{Summe} & & \textbf{218} & \\ \bottomrule
\end{tabular}
\caption{Automatisierte Tests (Stand 27.09.2026)}
\end{reporttable}

\subsubsection*{Golden-File-Tests}
Das Verzeichnis \code{interpreter/src/test/resources/programs} enthält neun thematische Programme (Grundlagen,
Kontrollfluss, Funktionen, Klassen, Vererbung, Zeiger, Scoping, \code{void main}, Exit-Codes), jeweils mit einer
\code{.expected}-Datei der erwarteten Ausgabe. Ein Kommentar \code{// expect-exit: N} legt einen abweichenden
Exit-Code fest. Dieselben Dateien dienen auch dem GCC-Vergleich.

\subsection{Differenzieller Vergleich mit GCC}\label{sec:gcc}
Die stärkste Korrektheitsaussage liefert der Vergleich mit einem echten C++-Compiler. Das Skript
\code{scripts/compare-gcc.sh} führt jedes Testprogramm (1) im Interpreter aus, (2) übersetzt es mit
\code{minicpp to-cpp} nach C++, kompiliert es mit \code{g++ -std=c++17 -fwrapv -O1} und führt es aus, und vergleicht
(3) beide Ausgaben und Exit-Codes mit der erwarteten Ausgabe. \code{-fwrapv} stellt sicher, dass GCC
Ganzzahlüberlauf ebenfalls als Zweierkomplement behandelt.

\subsubsection*{Continuous Integration}
Die GitHub-Actions-Pipeline (\code{.github/workflows/ci.yml}) läuft bei jedem Push und Pull-Request auf
\code{ubuntu-latest}: Build und Tests aller Module mit Temurin 21, Build der VS-Code-Extension mit Node 22 und
abschließend der GCC-Vergleich.

\begin{todo}
Aktuelles Ergebnis des GCC-Vergleichs aus der CI (Anzahl PASS/FAIL) eintragen; lokal ist kein g++ installiert.
\end{todo}

\subsection{Performance und Skalierbarkeit}\label{sec:perf}
Die Benchmarks im Modul \code{benchmark} verwenden JMH und messen die mittlere Laufzeit.
\code{InterpreterBenchmark} misst reine Interpretation (übersetzt wird einmal pro Trial) für vier Lasttypen mit
Größenparameter $n$; \code{FrontendBenchmark} misst Parsen bzw.\ Parsen und Prüfen generierter Programme mit 10, 100
und 1\,000 Klassen.

% Messwerte aus docs/jmh.json (ms/op)
\begin{figure}[htbp]
\centering
\begin{tikzpicture}
\begin{axis}[
  width=\linewidth, height=6.4cm,
  ybar, bar width=10pt,
  ymode=log, log origin=infty, ymin=0.1, ymax=200,
  symbolic x coords={fib,loop,vd,oc}, xtick=data,
  xticklabels={{Rekursion\\(fib($n$))}, {Schleife\\($n\cdot 10\,000$ Iter.)},
               {Virt.\ Dispatch\\($n\cdot 1\,000$ Aufrufe)}, {Objektkopien\\($n\cdot 1\,000$)}},
  xticklabel style={align=center, font=\small},
  enlarge x limits=0.15,
  ylabel={mittlere Laufzeit [ms], log}, ylabel style={color=ink2},
  ymajorgrids, grid style={gridc},
  axis lines*=left, axis line style={ink2}, tick style={ink2},
  log ticks with fixed point, /pgf/number format/use comma,
  legend style={draw=none, at={(0.02,0.98)}, anchor=north west, legend columns=3, font=\small},
  point meta=rawy,
  nodes near coords={\pgfmathprintnumber[fixed,precision=2,use comma]{\pgfplotspointmeta}},
  nodes near coords style={font=\tiny, color=ink2},
]
\addplot[fill=seqA, draw=none] coordinates {(fib,0.2293) (loop,6.1577) (vd,1.4993) (oc,9.1262)};
\addplot[fill=seqB, draw=none] coordinates {(fib,1.0810) (loop,8.1726) (vd,1.9995) (oc,13.2311)};
\addplot[fill=seqC, draw=none] coordinates {(fib,10.3365) (loop,10.0203) (vd,2.4108) (oc,16.1949)};
\legend{$n = 15$, $n = 20$, $n = 25$}
\end{axis}
\end{tikzpicture}
\caption{Laufzeit der Interpreter-Benchmarks nach Lasttyp und Größe $n$ (logarithmische Achse)}
\end{figure}

Die Rekursion \code{fib(n)} wächst erwartungsgemäß exponentiell (Faktor $\approx 1{,}6$ pro Schritt von $n$): von
0,23\,ms bei $n = 15$ auf 10,3\,ms bei $n = 25$, bei rund 243\,000 Aufrufen also etwa 40\,ns pro interpretiertem
Aufruf. Die drei anderen Lasten wachsen linear in $n$; sie zeigen jeweils einen konstanten Sockel, der auf
Thread-Start (512-MiB-Stack) und JIT-Aufwärmeffekte zurückgeht. Virtueller Dispatch ist mit 2,41\,ms für 25\,000
Aufrufe günstig, da vtables pro Klasse nur einmal aufgebaut werden. Objektkopien sind die teuerste Lastart, weil
jede Kopie neue Cells und Maps alloziert.

\begin{figure}[htbp]
\centering
\begin{tikzpicture}
\begin{axis}[
  width=\linewidth, height=5.8cm,
  xmode=log, ymode=log,
  xtick={10,100,1000}, xticklabels={10,100,1\,000},
  xmin=7, xmax=1500, ymin=0.08, ymax=200,
  xlabel={Anzahl generierter Klassen}, ylabel={mittlere Laufzeit [ms], log},
  label style={color=ink2},
  grid=major, grid style={gridc},
  axis lines*=left, axis line style={ink2}, tick style={ink2},
  log ticks with fixed point, /pgf/number format/use comma,
  legend style={draw=none, at={(0.02,0.98)}, anchor=north west, font=\small},
  point meta=rawy,
  nodes near coords={\pgfmathprintnumber[fixed,precision=2,use comma]{\pgfplotspointmeta}},
  nodes near coords style={font=\tiny, color=ink2},
]
\addplot[color=accent, line width=1.5pt, mark=*, mark options={fill=accent, draw=white},
         nodes near coords style={anchor=north, yshift=-3pt}]
  coordinates {(10,0.2157) (100,2.1472) (1000,21.3779)};
\addplot[color=orange, line width=1.5pt, mark=*, mark options={fill=orange, draw=white},
         nodes near coords style={anchor=south, yshift=3pt}]
  coordinates {(10,0.2767) (100,2.7404) (1000,61.1965)};
\legend{Parsen (Lexer, Parser, AST), Parsen + Resolver + Typprüfung}
\end{axis}
\end{tikzpicture}
\caption{Skalierung von Parsen und statischer Analyse mit der Programmgröße (log-log)}
\end{figure}

Das Frontend skaliert zwischen 10 und 100 Klassen linear (Faktor 10 bei zehnfacher Größe). Für 1\,000 Klassen
benötigt das Parsen 21,4\,ms, Parsen und Prüfen 61,2\,ms; der überproportionale Anstieg der Prüfung in diesem Punkt
ist auffällig, hat aber in diesem kurzen Messlauf auch die größte Streuung ($\pm$\,47,5\,ms) und muss mit längeren
Läufen bestätigt werden.

Messumgebung: JMH mit 1 Fork, $2 \times 1$\,s Warm-up und $3 \times 1$\,s Messung; AMD Ryzen 7 9800X3D, OpenJDK
64-Bit Server VM 25.0.1+8-LTS. Die Rohdaten liegen in \code{benchmark/build/jmh-result.json}.

\begin{todo}
Die Werte stammen aus einem kurzen Vorab-Lauf mit großer Streuung. Für die Endfassung mit den Standardeinstellungen
oder mehr Forks messen (\code{./gradlew :benchmark:jmh -Pjmh=\dq -f 3 -wi 5 -i 10\dq}) und ggf.\ einen Vergleich
mit \code{g++ -O0} ergänzen.
\end{todo}

\subsection{Bewertung}
Die Tests decken jede Phase einzeln sowie das Zusammenspiel über Golden Files ab; der GCC-Vergleich belegt, dass die
Semantik nicht nur mit den eigenen Erwartungen, sondern mit der Referenzimplementierung übereinstimmt. Besonders
wertvoll war die Kombination aus \code{ErrorTest} und Laufzeit-Lebensdauerprüfung: Jeder Fall von undefiniertem
Verhalten in C++ ist als eigener Testfall mit erwarteter Meldung festgehalten. Die Performance ist für den
Einsatzzweck -- Lehre, interaktive Nutzung, Editor-Analyse -- ausreichend: typische Beispielprogramme laufen im
einstelligen Millisekundenbereich, und die statische Analyse einer Datei mit 100 Klassen bleibt deutlich unter der
200-ms-Entprellzeit des Language-Servers.

% ================================================================== 8 Organisation
\section{Projektorganisation}\label{sec:organisation}
\subsection{Arbeitsplan}
Das Projekt war auf zwölf Wochen angelegt. Die Kernphasen wurden gemeinsam bearbeitet, die Erweiterungen (LSP und
GenAI) parallel von jeweils einem Teammitglied.

\begin{reporttable}
\begin{tabular}{L{0.1}L{0.68}L{0.14}}
\thead \hcell{Woche} & \hcell{Schwerpunkt} & \hcell{Verantw.} \\ \midrule
1 & Repository (Gradle), ANTLR4 einrichten, Grammatik-Entwurf, CI & C + D \\
2--4 & ANTLR4-Lexer und -Parser, AST-Klassenhierarchie, Fehlerbehandlung & C + D \\
5--6 & Two-Pass-Resolver, Typprüfung, semantische Prüfungen & C + D \\
7--8 & Tree-Walking-Interpreter, REPL, Objektmodell, vtables & C + D \\
9--11 & LSP4J-Server, VS-Code-Extension & C \\
9--11 & GenAI-Server, Ollama-Anbindung, Mock-Provider & D \\
12 & Benchmarks, Korrektheitstests, Dokumentation, Walk-Through & C + D \\ \bottomrule
\end{tabular}
\caption{Arbeitsplan (C = Clemens Vogtländer, D = Dennis Gorpinic)}
\end{reporttable}

\subsection{Aufgabenteilung}
\begin{reporttable}
\begin{tabular}{L{0.56}L{0.17}L{0.17}}
\thead \hcell{Aufgabe} & \hcell{C. Vogtländer} & \hcell{D. Gorpinic} \\ \midrule
Grammatik (Lexer und Parser), AST-Hierarchie & 50\,\% & 50\,\% \\
Resolver, Typprüfung und Überladungsauflösung & 50\,\% & 50\,\% \\
Interpreter, Objektmodell (vtable, Slicing), REPL & 50\,\% & 50\,\% \\
LSP-Server und VS-Code-Extension & 100\,\% & -- \\
GenAI-Server, Ollama-Anbindung, Prompt-Engineering & -- & 100\,\% \\
Mock-Provider für Tests & 20\,\% & 80\,\% \\
Korrektheitstests, GCC-Vergleich, JMH-Benchmarks & 50\,\% & 50\,\% \\
Abschlussdokumentation und Walk-Through & 50\,\% & 50\,\% \\ \bottomrule
\end{tabular}
\caption{Aufgabenteilung}
\end{reporttable}

Entwickelt wurde in einem gemeinsamen Git-Repository; die Implementierung des Interpreters liegt auf dem Branch
\code{interpreter}. Die CI prüft jeden Push, sodass Regressionen im gemeinsam genutzten Kern früh sichtbar werden --
wichtig, weil LSP- und GenAI-Server direkt von dessen AST-Annotationen abhängen.

\begin{todo}
Tatsächlichen Verlauf gegenüber dem Plan reflektieren: Was hat länger gedauert (z.\,B. Zeiger-Semantik, die
nachträglich in den Umfang aufgenommen wurde), welche Entscheidungen wurden revidiert, was würde das Team anders
machen?
\end{todo}

% ================================================================== 9 Ausblick
\section{Grenzen und Ausblick}\label{sec:ausblick}
\subsection{Bekannte Grenzen}
\begin{itemize}
  \item \textbf{Kein AST bei Syntaxfehlern.} Die Pipeline bricht nach Syntaxfehlern ab. Der Language-Server
    überbrückt dies für die Completion mit der letzten parsebaren Version; Hover und Navigation sind in einer
    syntaktisch fehlerhaften Datei jedoch nicht verfügbar. Eine Fehlerrecovery mit \code{ErrorExpr}-Knoten würde
    partielle ASTs ermöglichen.
  \item \textbf{Erste Fehlerphase gewinnt.} Meldet der Resolver Fehler, entfällt die Typprüfung; Typfehler werden
    erst nach Behebung der Namensfehler sichtbar.
  \item \textbf{Vollständige Neuanalyse.} Jede Änderung analysiert das gesamte Dokument neu (siehe
    Abschnitt~\ref{sec:sync}).
  \item \textbf{Interpretationsgeschwindigkeit.} Der Tree-Walker schlägt Variablen über eine
    \code{IdentityHashMap} pro Frame nach und alloziert für jeden Wert eine Cell. Für den Lehrkontext genügt das;
    für rechenintensive Programme wäre er um Größenordnungen langsamer als kompilierter Code.
  \item \textbf{Ein Dokument pro Programm.} MiniC++ kennt keine Module; der Language-Server behandelt jede Datei
    unabhängig.
\end{itemize}

\subsection{Mögliche Weiterentwicklungen}
\begin{itemize}
  \item \textbf{Slot-Auflösung:} Der Resolver könnte jeder lokalen Variable einen Index im Frame zuweisen, sodass
    Zugriffe zu Array-Operationen werden -- der naheliegendste Performancegewinn.
  \item \textbf{Bytecode oder Truffle:} Eine Übersetzung in einen kompakten Bytecode oder eine Implementierung auf
    GraalVM Truffle würde JIT-Kompilierung ermöglichen.
  \item \textbf{Sprachumfang:} \code{for}, \code{break}/\code{continue}, \code{const} und Arrays sind die am
    häufigsten vermissten Konstrukte und ließen sich in die bestehende Architektur einfügen.
  \item \textbf{Debug Adapter Protocol:} Da der Interpreter jeden Knoten mit Quellbereich ausführt, liegt ein
    Debugger für VS Code (Breakpoints, Variablenansicht) nahe.
  \item \textbf{MCP-Anbindung:} Die Assistenzfunktionen zusätzlich als Model-Context-Protocol-Tools bereitstellen,
    sodass KI-Agenten MiniC++-Code prüfen und ausführen können.
\end{itemize}

\subsection{Fazit}
Das Projekt zeigt, dass eine klar abgegrenzte Sprache mit einem sauber geschichteten Frontend eine tragfähige
Grundlage für ein ganzes Werkzeug-Ökosystem ist. Die Entscheidung, alle Analyseergebnisse als Annotationen am AST und
alle Probleme als positionsgenaue \code{Diagnostic}s zu modellieren, hat sich ausgezahlt: Interpreter, REPL,
Language-Server und GenAI-Server nutzen dieselbe Pipeline und liefern konsistente Ergebnisse. Die Korrektheit
gegenüber C++ ist durch den differenziellen GCC-Vergleich belegt, und die Laufzeitprüfungen machen undefiniertes
Verhalten, das in C++ schwer zu finden ist, für Lernende sichtbar.

\begin{todo}
Fazit nach Abschluss der Evaluation überarbeiten; ggf.\ Abschnitt zu Lessons Learned ergänzen.
\end{todo}

% ================================================================== Anhang
\clearpage
\appendix
\section{Build und Nutzung}\label{app:build}
Voraussetzung ist ein JDK 21 oder neuer; für die Extension zusätzlich Node.js.

\begin{lstlisting}[language=bash,caption={Build-Befehle}]
./gradlew build                          # alle Module bauen und testen
./gradlew :interpreter:installDist       # -> interpreter/build/install/minicpp/bin/minicpp
./gradlew :lsp:installDist               # -> lsp/build/install/minicpp-lsp/bin/minicpp-lsp
./gradlew :mcp:run --args="--provider mock"     # GenAI-Server ohne Ollama auf :8080
./gradlew :benchmark:jmh                 # Benchmarks, Ergebnis in benchmark/build/jmh-result.json
cd vscode && npm install && npm run package      # -> minicpp-1.0.0.vsix
\end{lstlisting}

\begin{reporttable}
\begin{tabular}{L{0.33}L{0.59}}
\thead \hcell{Befehl} & \hcell{Wirkung} \\ \midrule
\code{minicpp run datei.cpp} & Programm ausführen; Exit-Code = Rückgabewert von \code{main} \\
\code{minicpp check datei.cpp} & nur Syntax-, Namens- und Typprüfung \\
\code{minicpp repl [datei.cpp]} & REPL, optional mit vorher geladener Datei \\
\code{minicpp ast datei.cpp} & AST ausgeben \\
\code{minicpp to-cpp datei.cpp} & nach Standard-C++17 übersetzen \\ \bottomrule
\end{tabular}
\caption{Unterbefehle der Kommandozeile}
\end{reporttable}

\begin{reporttable}
\begin{tabular}{L{0.28}L{0.32}L{0.32}}
\thead \hcell{Option} & \hcell{Umgebungsvariable} & \hcell{Standard} \\ \midrule
\code{--host} & \code{MCP\_HOST} & \code{127.0.0.1} \\
\code{--port} & \code{MCP\_PORT} & \code{8080} \\
\code{--provider} & \code{MCP\_PROVIDER} & \code{ollama} (alternativ \code{mock}) \\
\code{--ollama-url} & \code{OLLAMA\_URL} & lokale Ollama-Instanz \\
\code{--model} & \code{OLLAMA\_MODEL} & \code{codellama} \\
\code{--timeout} & \code{OLLAMA\_TIMEOUT} & 120\,s \\ \bottomrule
\end{tabular}
\caption{Konfiguration des GenAI-Servers}
\end{reporttable}

\section{Grammatik: Deklarationen und Anweisungen}
\begin{lstlisting}[caption={Parser-Regeln für Deklarationen und Anweisungen (vollständig in \code{MiniCpp.g4})}]
program       : declaration* EOF ;
replInput     : (declaration | statement)* expr? EOF ;
declaration   : functionDef | classDef ;
classDef      : CLASS Identifier (COLON PUBLIC Identifier)?
                LBRACE PUBLIC COLON memberDecl* RBRACE SEMI ;
memberDecl    : fieldDecl | constructorDef | methodDef ;
fieldDecl     : type Identifier SEMI ;
constructorDef: Identifier LPAREN paramList? RPAREN block ;
methodDef     : VIRTUAL? type Identifier LPAREN paramList? RPAREN block ;
functionDef   : type Identifier LPAREN paramList? RPAREN block ;
param         : (type | typeRef) Identifier ;
typeRef       : type AMP ;
varDecl       : type Identifier
              | type Identifier LPAREN argList RPAREN
              | (type | typeRef) Identifier ASSIGN expr ;
ifStmt        : IF LPAREN expr RPAREN statement (ELSE statement)? ;
whileStmt     : WHILE LPAREN expr RPAREN statement ;
returnStmt    : RETURN expr? ;
deleteStmt    : DELETE expr ;
newExpr       : NEW (baseType | Identifier) (LPAREN argList? RPAREN)? ;
\end{lstlisting}

\section{Abkürzungen}
\begin{reporttable}
\begin{tabular}{L{0.2}L{0.72}}
\thead \hcell{Abkürzung} & \hcell{Bedeutung} \\ \midrule
ADR & Architecture Decision Record \\
AST & Abstract Syntax Tree, abstrakter Syntaxbaum \\
CI & Continuous Integration \\
JMH & Java Microbenchmark Harness \\
LSP & Language Server Protocol \\
MCP & Model Context Protocol \\
REPL & Read-Eval-Print-Loop \\
vtable & Tabelle virtueller Methoden einer Klasse \\ \bottomrule
\end{tabular}
\end{reporttable}

\end{document}
